% Pursuing Stacks (cotes 134-2 to 134-8): archive cote 134-8, site page 20 (PDF page 21), typescript page 574.
% typescript. Draft compared: ps7.tex:921-985.
% AI-transcribed by Claude Opus 5.5 (claude-opus-5-5, adaptive thinking, effort high), Message Batches API, batch ps4 (msgbatch_012Aoc1NuskATPsyBeKSYK7J), 2026-11-12; not yet checked by eye.
\scanpage{134-8/20}
\mspage{574}

\uline{of} (\uline{locally compac}t) \uline{topological spaces}, as worked out by Verdier in one
of his Bourbaki talks - it being understood that when applying the
latter formalism to spaces endowed with conical stratifications, maps
between these which are compatible with the stratifications (in a suitab%
% NOTE: the right edge of the scan cuts off the ends of some lines (after "suitab", "satis", "one"); a hyphen or letter may be hidden there
le sense which should still be pinned down), and to sheaves of modules
which \xout{satisfy the} are compatible with the stratifications in the sense
above, these will give rise (via the ``six operations'') to sheaves satis
-fying the same compatibility. (NB when speaking of``sheaves'', I really
mean complexes of sheaves,\inserted{$K^\bullet$} and the compatibility condition should be
understood as a condition \deleted{at}\inserted{for} the ordinary sheaves of modules $\underline{H}^i(K^\bullet)$.)
This remark should allow to work out quite explicitly, in purely alge-%
braic (or ``combinatorial'') terms, the ``six operations'' in the context
of locally finite ordered sets, at any rate.

This interpretation suggests that an ordered set I should be viewed
as a ``proper'' object of (Cat) iff I is finite. In the same vein, when-%
ever I is ``locally proper'' namely locally finite, and moreover it's to-%
pological realization I is a ``smooth'' topological space in the usual
sense, namely is a topological \uline{variety} (for which it is enough that
for every x in I, the topological realization
\hand{(7)}
\[ \hand{|}I_{\hand{>}x}\hand{|} \]
of the set of elements y such that y \hand{$>$} x should be a sphere), we would
consider I as a ``smooth'' object in (Cat). If A is any object in (Cat)
and I is an ordered set which is finite resp. ``smooth'' in the sense
above, we will surely expect AxI \hand{$\longrightarrow$} A to be``proper'' resp. ``smooth''
for the duality formalism we wish to develop in (Cat).

It should be kept in mind that for the \inserted{algebraic} interpretation above\inserted{to hold}, for
sheaves on a conically stratified space \hand{$|$}I\hand{$|$} in terms of an ordered set
I , we had to take on the indexing set\inserted{I} for the strata the order relati-%
on \uline{opposite} to the inclusion relation between (closed) strata - other-%
wise the correct interpretation of sheaves constant on the open strata
is via \uline{covariant} functors I \hand{$\longrightarrow$} (Sets), i.e. \uline{presheaves on} $I^o$ (not I).
At any rate, as there is a canonical homeomorphism
\hand{(8)}
\[ \hand{|}I\hand{|} \hand{\simeq} \hand{|}I^o\hand{|} \]
respecting\inserted{the canonical} barycentric subdivisions of both sides, notions \xout{\unsure{f}}for I
(such as properness, or smoothness) which are expressed as intrinsic \xout{p}
properties of the corresponding topological space \hand{$|$}I\hand{$|$} (independently of
it's subdivision) are autodual - they hold for I iff they do for $I^o$.

\margin{20}
